\dykband{EXERCISES IN SWAFL}
{\small Open SWAFL (desktop or online) with the benchmark project. The chapter of the \href{https://github.com/halilunaziru73-creator/Coupled-Water-Fluxes-Pore-Scale-Dynamics-Precision-Irrigation-Mulching-Feedback-Loops/raw/main/Software/v1.1.0/SWAFL_Desktop/docs/manual_v1.1.0/SWAFL_Manual_v1.1.0.pdf}{SWAFL manual} that explains each task is given in brackets; record your values in a table and compare them with what is expected (\emph{Check}).}\par\vspace{3pt}
\begin{exlist}
\item {\color{khblue}[Manual, Chapter 25]} Open the Daisy workbench and the process map. Click the boxes of root water uptake and soil water and read their descriptions. List the inputs that each box needs.\par\noindent{\small\itshape Check: uptake needs the root density, the suction and the potential transpiration.}
\item {\color{khblue}[Manual, Chapter 28]} In \emph{Experiments and scenarios} choose an irrigation scenario for the loam and look at the plan of the runs before starting. How many runs does it hold, and which factors are varied?\par\noindent{\small\itshape Check: the plan lists every run with its factors and the budget.}
\item {\color{khblue}[Manual, Chapter 30]} After a run, open \emph{Results and analysis} and the plot of soil water suction (pF) over depth and time. Identify a dry summer and the wetting by irrigation.\par\noindent{\small\itshape Check: the pF rises in summer in the upper layers and falls after each irrigation.}
\item {\color{khblue}[Manual, Chapter 7]} In \emph{Profile evolution} compare the profile of the loam and of the sand two days after the same rain. Explain the difference with the Richards equation.\par\noindent{\small\itshape Check: the sand drains faster and holds less water near the surface.}
\end{exlist}
